%% =============================================================================
%% ABSTRACT (English + French + Arabic)
%% =============================================================================

\clearpage
\phantomsection
\addcontentsline{toc}{chapter}{Abstract}

\begin{center}
\textbf{\Large Abstract}
\end{center}
\vspace{0.2cm}
\indent Colorectal cancer is a massive global problem, causing roughly 10\% of all cancer cases. The most effective intervention is the early detection and resection of suspicious polyps. However, navigating the 2,000 cm\textsuperscript{2} surface area of the colon with a flexible endoscope presents a significant technical challenge. It generally requires two coordinated physicians, relying on rapid decision-making and steady motor control under pressure. Rather than relying on costly robotic hardware, we pursued a software-centric approach. We developed a semi-autonomous agent to manage the core navigation tasks within a high-fidelity virtual colon. To manage complexity, we divided the challenge: one agent steers the endoscope tip through the colonic lumen, and once a polyp is identified, a second agent takes over to accurately position a wire snare around it. Both utilize Deep Reinforcement Learning with Proximal Policy Optimization (PPO), Convolutional Neural Networks (CNN), and a Long Short-Term Memory (LSTM) core for spatial awareness. Consequently, polypectomy extends beyond basic image recognition; it is a highly complex, sequential decision-making problem. By structuring the reward system to incentivize safe maneuvers and penalize collisions, the agents learned to balance the preservation of healthy tissue with the execution of precise resections. Our unique dual-sphere proxy proved that these vision-only agents can navigate and capture polyps under realistic, noisy conditions.

\vspace{0.5cm}
\noindent\textbf{Keywords:} colorectal cancer, deep reinforcement learning, PPO, CNN, virtual simulation, polypectomy, snare capture

\clearpage
\begin{center}
\textbf{\Large R\'esum\'e}
\end{center}
\vspace{0.2cm}
\indent Le cancer colorectal cause environ 10 \% des cas de cancer mondiaux. La meilleure solution reste la détection et la résection précoce des polypes. Cependant, naviguer dans un côlon de 2 000 cm\textsuperscript{2} avec un endoscope délicat est un défi technique majeur. Cette procédure nécessite généralement deux médecins coordonnés, reposant sur leur précision et rapidité sous une pression intense. Plutôt que de développer des robots coûteux, nous avons privilégié une approche logicielle. Nous avons créé un agent semi-autonome pour accomplir le gros du travail dans un environnement virtuel hyper-réaliste. Pour gérer la complexité, nous avons divisé le défi : un agent dirige l'embout de l'endoscope à travers les tunnels roses (côlon), et une fois un polype repéré, un deuxième agent prend le relais pour refermer délicatement une anse métallique autour de lui. Les deux utilisent l'Apprentissage par Renforcement Profond avec Proximal Policy Optimization (PPO), des Réseaux de Neurones Convolutifs (CNN) et une mémoire Long Short-Term Memory (LSTM) pour ne pas se perdre. La polypectomie s'avère être un problème décisionnel séquentiel complexe. En entraînant nos agents par renforcement, ils ont appris à équilibrer la précision de la coupe avec la préservation des tissus sains. Notre modèle montre que des agents basés uniquement sur la vision peuvent naviguer et réussir des interventions sous des conditions de bruit réalistes.

\vspace{0.5cm}
\noindent\textbf{Mots-cl\'es:} cancer colorectal, apprentissage par renforcement profond, PPO, CNN, simulation virtuelle, polypectomie, capture par anse

\clearpage
\begin{Arabic}
\begin{center}
\textbf{\Large ملخص}
\end{center}
\vspace{0.2cm}
\indent يُعد سرطان القولون والمستقيم صداعاً عالمياً هائلاً، حيث يتسبب في حوالي \textenglish{10\%} من جميع حالات السرطان. أفضل حل؟ اكتشاف الأورام الحميدة المشبوهة وقطعها مبكراً. لكن توجيه منظار دقيق داخل قولون بمساحة \textenglish{2000} سم\textsuperscript{2} أمر بالغ الصعوبة! يتطلب الأمر عادةً طبيبين منسقين، حيث يعتمد كلياً على قراراتهما السريعة وأيديهما الثابتة تحت ضغط شديد. بدلاً من بناء روبوتات باهظة الثمن، اتخذنا نهجاً برمجياً جريئاً! قمنا بإنشاء وكيل شبه مستقل للقيام بالعمل الشاق داخل قولون افتراضي فائق الواقعية. ولإبقاء الأمور تحت السيطرة، قمنا بتقسيم التحدي: حيث يقوم الوكيل الأول بتوجيه طرف المنظار عبر الأنفاق الوردية (القولون)، وبمجرد اكتشاف ورم، يتولى وكيل ثانٍ مهمة إغلاق حلقة سلكية برفق حوله. كلاهما يستخدم التعلم المعزز العميق (\textenglish{PPO}) والشبكات العصبية التلافيفية (\textenglish{CNN}) مع ذاكرة (\textenglish{LSTM}) حتى لا يضيعا. اتضح أن استئصال السليلة ليس مجرد تعرف أساسي على الصور؛ بل هو لعبة قرارات متسلسلة في غاية التعقيد! من خلال إعطاء الذكاء الاصطناعي الخاص بنا مكافآت على الحركات الجيدة وعقوبات عند الاصطدام، تعلموا الموازنة بشكل مثالي بين الحفاظ على الأنسجة السليمة وإجراء عمليات قطع دقيقة. أثبت نظام الوكيل مزدوج الكرات الفريد لدينا أن هؤلاء الوكلاء المعتمدين على الرؤية فقط يمكنهم التنقل والتقاط الأورام الحميدة في ظل ظروف واقعية صاخبة.

\vspace{0.5cm}
\noindent\textbf{الكلمات المفتاحية:} سرطان القولون والمستقيم، التعلم المعزز العميق، \textenglish{PPO}، \textenglish{CNN}، محاكاة افتراضية، استئصال السليلة، الالتقاط بالحلقة السلكية.
\end{Arabic}

\clearpage